% Modul Belajar 15, dihasilkan oleh tools/build.py dari tools/konten/p15.py
\documentclass[a4paper]{article}
\usepackage{modulITERA}
\begin{document}
\Sampul{15}{Review dan Simulasi UAS}{Merangkai materi Pertemuan 9 sampai 14 dan berlatih dengan soal bergaya UAS}{CPMK0607 (menerapkan) dan CPMK0608 (menjelaskan) pada konsep dasar algoritma dan sistem komputer}{sekitar 3 jam belajar mandiri, ditambah 150 menit di kelas}{\texttt{02 Hands-on/P15 - Review UAS - Hands-on/}}
\Bagian{Modul 15}{Peta belajar}
Ikuti urutan ini dari atas ke bawah. Setiap bagian memakai apa yang sudah dibahas di bagian sebelumnya.
\par\vspace{4pt}\noindent{\fontsize{9.5pt}{13pt}\selectfont
\begin{tabular}{@{}>{\raggedright\arraybackslash}p{0.140\textwidth}>{\raggedright\arraybackslash}p{0.840\textwidth}@{}}
\kepala{Urutan} & \kepala{Bagian}\\
1 & Tentang modul ini\\\garis
2 & Masalah pembuka: Satu soal, banyak materi\\\garis
3 & Langkah 1: Peta materi UAS\\\garis
4 & Langkah 2: Strategi soal A\\\garis
5 & Langkah 3: Strategi soal B\\\garis
6 & Langkah 4: Contoh tracing gabungan\\\garis
7 & Pesan error yang sering muncul\\\garis
8 & Latihan mandiri\\\garis
9 & Rangkuman\\\garis
10 & Cek pemahaman\\\garis
11 & Exit Ticket\\\garis
12 & Bacaan lanjut dan persiapan minggu depan\\
\garisdasar
\end{tabular}}\par\vspace{6pt}
\Bagian{Pembuka}{Tentang modul ini}
Modul ini mendampingi Pertemuan 15. Tidak ada materi baru. Tujuannya merangkai materi Pertemuan 9 sampai 14 dan berlatih dengan soal bergaya UAS: soal A menulis program, soal B menelusuri dan menjelaskan.

Isinya sama dengan slide di kelas, tetapi ditulis lebih pelan dan lebih lengkap, supaya kalian bisa mengulang sendiri di rumah tanpa harus menebak maksud slide.

\Sub{Setelah menyelesaikan modul ini, kalian bisa}
\par\vspace{4pt}\noindent{\fontsize{9.5pt}{13pt}\selectfont
\begin{tabular}{@{}>{\raggedright\arraybackslash}p{0.070\textwidth}>{\raggedright\arraybackslash}p{0.680\textwidth}>{\raggedright\arraybackslash}p{0.230\textwidth}@{}}
\kepala{No} & \kepala{Kemampuan} & \kepala{Dibahas di}\\
1 & Menyelesaikan masalah yang menggabungkan array, string, fungsi, pengurutan atau pencarian, dan struct dalam satu program C++ dalam batas waktu ujian (Sub-CPMK 15.1, CPMK0607) & Langkah 1 sampai 4\\\garis
2 & Menelusuri program yang memakai array, fungsi, dan rekursi, lalu menjelaskan setiap perubahan nilai dengan istilah yang tepat (Sub-CPMK 15.2, CPMK0608) & kotak Tebak dulu, Latihan 3\\
\garisdasar
\end{tabular}}\par\vspace{6pt}
Karena setiap instrumen penilaian dibagi rata antara Bagian A (menulis kode) dan Bagian B (menelusuri dan menjelaskan), yang dinilai bukan hanya program kalian jalan, tetapi juga kalian bisa menjelaskan mengapa program itu jalan.

\Sub{Cara memakai modul ini}
Setiap langkah punya pola yang sama: penjelasan konsep, kadang contoh dari kehidupan sehari-hari, lalu kode yang bisa langsung kalian ketik. Sesudah kode ada kotak \textbf{Tebak dulu}. Tulis tebakan kalian di kertas sebelum membaca \textbf{Jawaban} di bawahnya. Kebiasaan menebak lalu membuktikan inilah yang membuat konsep menempel, bukan sekadar membaca.

\begin{Penting}[Ketik, jangan salin]
Ketik ulang setiap contoh di GDB Online atau editor kalian, jangan disalin-tempel. Saat mengetik, kalian akan membuat salah ketik, membaca pesan error, dan memperbaikinya. Itu bagian dari belajar.
\end{Penting}
\Sub{Yang perlu disiapkan}
Peramban dengan akses ke onlinegdb.com, atau g++ di komputer kalian sendiri. Kertas dan alat tulis untuk menebak keluaran dan membuat trace table. Folder hands-on pertemuan ini dari repo mata kuliah.

\Bagian{Pembuka}{Masalah pembuka: Satu soal, banyak materi}
Soal UAS biasanya menggabungkan beberapa materi sekaligus. Contohnya ranking peserta: data disimpan sebagai array of struct, diurutkan dengan fungsi, lalu dicari namanya.

Hari ini kita berlatih memecah soal seperti itu menjadi bagian-bagian yang sudah kalian kuasai.

\begin{MKeluaran}[title={Tampilan yang diinginkan}]
1. Siti 95
2. Andi 88
3. Rani 82
4. Budi 75
5. Doni 70
Andi peringkat 2
\end{MKeluaran}
\begin{Tebak}[Coba pikirkan]
Materi pertemuan mana saja yang dipakai soal ini? Bagian mana yang menurut kalian paling rawan salah?
\end{Tebak}
\begin{Jawaban}[Cara memecahnya]
Baca soal: Garis bawahi masukan, keluaran, dan aturan. Pecah: Petakan setiap bagian ke struct, fungsi, sort, atau search. Uji: Siapkan data uji, termasuk data yang tidak ditemukan dan nilai sama.
\end{Jawaban}
\Bagian{Langkah 1}{Peta materi UAS}
Tabel ini merangkum konsep kunci dan jebakan yang paling sering muncul di setiap pertemuan. Gunakan sebagai daftar periksa: untuk setiap baris, pastikan kalian bisa menulis contoh kode dan menjelaskan jebakannya.

\par\vspace{4pt}\noindent{\fontsize{9.5pt}{13pt}\selectfont
\begin{tabular}{@{}>{\raggedright\arraybackslash}p{0.142\textwidth}>{\raggedright\arraybackslash}p{0.392\textwidth}>{\raggedright\arraybackslash}p{0.446\textwidth}@{}}
\kepala{Pertemuan} & \kepala{Konsep kunci} & \kepala{Jebakan yang sering muncul}\\
P9 & array 1D, maksimum, pergeseran & indeks di luar batas, \texttt{i <= n}\\\garis
P10 & array 2D, \texttt{string}, \texttt{cctype} & urutan indeks [b][k], \texttt{'A'} vs \texttt{"A"}\\\garis
P11 & linear dan binary search & \texttt{low < high}, arah pergeseran terbalik\\\garis
P12 & bubble dan selection sort & tukar tanpa variabel sementara, kurang putaran\\\garis
P13 & fungsi, prototipe, referensi & pass by value pada fungsi tukar, jalur tanpa \texttt{return}\\\garis
P14 & rekursi, struct & basis salah, lupa \texttt{;} sesudah struct\\
\garisdasar
\end{tabular}}\par\vspace{6pt}
\Bagian{Langkah 2}{Strategi soal A}
Soal A di paruh kedua semester biasanya lebih panjang. Memecahnya menjadi fungsi bukan hanya gaya penulisan; itu cara paling aman agar sebagian program sudah pasti benar walaupun waktu habis.

\par\vspace{4pt}\noindent{\fontsize{9.5pt}{13pt}\selectfont
\begin{tabular}{@{}>{\raggedright\arraybackslash}p{0.280\textwidth}>{\raggedright\arraybackslash}p{0.700\textwidth}@{}}
\kepala{Bagian} & \kepala{Isi}\\
Rancang struktur & Tentukan struct dan daftar fungsi beserta parameternya sebelum mengetik.\\\garis
Bangun bertahap & Tulis satu fungsi, panggil dari \texttt{main}, kompilasi, uji. Lanjut ke fungsi berikutnya.\\\garis
Uji kasus tepi & Data kosong, satu data, nilai sama, dan data yang tidak ditemukan.\\
\garisdasar
\end{tabular}}\par\vspace{6pt}
\begin{Penting}[]
Fungsi yang sudah teruji satu per satu jauh lebih mudah dirangkai daripada satu \texttt{main} panjang yang belum pernah dijalankan.
\end{Penting}
\Bagian{Langkah 3}{Strategi soal B}
Soal B dinilai dari ketepatan langkah dan alasan. Untuk array, tuliskan isi array setelah setiap perubahan. Untuk rekursi, pohon pemanggilan atau tumpukan pemanggilan adalah alat utamanya.

\par\vspace{4pt}\noindent{\fontsize{9.5pt}{13pt}\selectfont
\begin{tabular}{@{}>{\raggedright\arraybackslash}p{0.280\textwidth}>{\raggedright\arraybackslash}p{0.700\textwidth}@{}}
\kepala{Bagian} & \kepala{Isi}\\
Array: tulis isinya & Tulis isi array sesudah setiap putaran atau setiap pemanggilan.\\\garis
Rekursi: gambar pohon & Tulis setiap pemanggilan beserta argumennya, lalu nilai kembali dari bawah ke atas.\\\garis
Jelaskan & Pakai istilah tepat: kasus basis, referensi, indeks batas, putaran.\\
\garisdasar
\end{tabular}}\par\vspace{6pt}
\begin{Penting}[]
Jawaban B tanpa langkah hanya mendapat sebagian poin, walaupun hasil akhirnya benar.
\end{Penting}
\Bagian{Langkah 4}{Contoh tracing gabungan}
Contoh ini menggabungkan array, fungsi, rekursi, dan operator ternary. Telusuri dari \texttt{hitung(a, 6, 4)} turun sampai \texttt{hitung(a, 0, 4)}, lalu jumlahkan nilai \texttt{c} saat kembali naik.

\begin{MKode}{gabungan.cpp}
int hitung(int a[], int n, int x) {
    if (n == 0) return 0;
    int c = (a[n - 1] > x) ? 1 : 0;
    return c + hitung(a, n - 1, x);
}

int main() {
    int a[6] = {7, 2, 9, 4, 9, 1};
    cout << hitung(a, 6, 4) << endl;
    return 0;
}
\end{MKode}
\begin{Tebak}[]
Sebelum menjalankan program, tulis keluarannya di kertas, karakter demi karakter.
\end{Tebak}
\begin{Jawaban}[]
\begin{MKeluaran}[]
3
\end{MKeluaran}
Fungsi menghitung elemen yang lebih besar dari x secara rekursif dari belakang: 9, 9, dan 7, jadi 3.
\end{Jawaban}
\Bagian{Waspada}{Pesan error yang sering muncul}
Semua pesan di tabel ini dihasilkan g++ dengan opsi \texttt{-Wall}, sama seperti yang dipakai \texttt{cek.sh}.
\par\vspace{4pt}\noindent{\fontsize{9.5pt}{13pt}\selectfont
\begin{tabular}{@{}>{\raggedright\arraybackslash}p{0.240\textwidth}>{\raggedright\arraybackslash}p{0.270\textwidth}>{\raggedright\arraybackslash}p{0.470\textwidth}@{}}
\kepala{Kesalahan} & \kepala{Contoh} & \kepala{Pesan pertama dari g++}\\
Lupa ; sesudah struct & \texttt{struct P \{ int x; \}} & \texttt{error: expected ';' after struct definition}\\\garis
Fungsi dipanggil sebelum dikenal & \texttt{panggil sebelum definisi} & \texttt{error: 'f' was not declared in this scope}\\\garis
Array 2D tanpa ukuran kolom & \texttt{void f(int m[][])} & \texttt{error: declaration of 'm' as multidimensional array must have bounds for all ...}\\\garis
Literal ke parameter referensi & \texttt{tukar(a, 5)} & \texttt{error: cannot bind non-const lvalue reference of type 'int\&' to an rvalue of ...}\\\garis
Tidak ada return di semua jalur & \texttt{if (...) return 1; saja} & \texttt{warning: control reaches end of non-void function}\\
\garisdasar
\end{tabular}}\par\vspace{6pt}
Kelima error ini mewakili kesalahan yang paling sering membuat program UAS tidak terkompilasi atau berperilaku acak. Biasakan mengompilasi setiap selesai menulis satu fungsi.

\Bagian{Latihan}{Latihan mandiri}
Latihan ini sama dengan hands-on di kelas. Semua file ada di \texttt{02 Hands-on/P15 - Review UAS - Hands-on/latihan/}. Kerjakan berurutan, lalu cocokkan dengan \texttt{expected/} atau jalankan \texttt{bash cek.sh}.
\par\vspace{4pt}\noindent{\fontsize{9.5pt}{13pt}\selectfont
\begin{tabular}{@{}>{\raggedright\arraybackslash}p{0.070\textwidth}>{\raggedright\arraybackslash}p{0.360\textwidth}>{\raggedright\arraybackslash}p{0.150\textwidth}>{\raggedright\arraybackslash}p{0.400\textwidth}@{}}
\kepala{No} & \kepala{File} & \kepala{Tingkat} & \kepala{Yang dilatih}\\
1 & \texttt{handson1-kata.cpp} & Dasar & array string, fungsi, karakter pertama\\\garis
2 & \texttt{handson2-ranking.cpp} & Menengah & struct, fungsi pengurutan, pencarian\\\garis
3 & \texttt{handson3-telusur.cpp} & Tantangan & pohon pemanggilan dan isi array\\\garis
+ & \texttt{bonus-matriks.cpp} & Pengayaan & fungsi dengan parameter array 2D\\
\garisdasar
\end{tabular}}\par\vspace{6pt}
\Sub{Latihan 1: handson1-kata.cpp}
Kata terpanjang lewat fungsi dan semua kata yang diawali vokal. Simulasi soal A.

\Sub{Latihan 2: handson2-ranking.cpp}
Ranking peserta dengan array of struct, fungsi selection sort, dan pencarian nama.

\Sub{Latihan 3: handson3-telusur.cpp}
Simulasi soal B: gambar pohon pemanggilan dan isi array, tulis perkiraan keluaran sebelum menjalankan. File ini sudah lulus \texttt{cek.sh}; yang dinilai adalah tracing kalian.

\Sub{Bonus: bonus-matriks.cpp}
Periksa apakah matriks simetris lewat fungsi \texttt{bool}, lalu tampilkan jumlah setiap baris.

\Sub{Latihan menelusuri}
\begin{MKode}{tebak.cpp}
struct P { int x; };

void naik(P p) { p.x++; }
void naikRef(P &p) { p.x++; }

int main() {
    P a = {5};
    naik(a);
    naikRef(a);
    naikRef(a);
    cout << a.x << endl;
    return 0;
}
\end{MKode}
\begin{Tebak}[]
Tulis keluarannya di kertas tanpa menjalankan program. Buat trace table bila perlu.
\end{Tebak}
\begin{Jawaban}[]
\begin{MKeluaran}[]
7
\end{MKeluaran}
\texttt{naik} menerima salinan sehingga \texttt{a.x} tidak berubah. Dua pemanggilan \texttt{naikRef} menambah \texttt{a.x} aslinya dua kali: 5 menjadi 7.
\end{Jawaban}
\Bagian{Penutup}{Rangkuman}
\par\vspace{4pt}\noindent{\fontsize{9.5pt}{13pt}\selectfont
\begin{tabular}{@{}>{\raggedright\arraybackslash}p{0.320\textwidth}>{\raggedright\arraybackslash}p{0.660\textwidth}@{}}
\kepala{Konsep} & \kepala{Yang perlu diingat}\\
Peta materi UAS & Tabel ini merangkum konsep kunci dan jebakan yang paling sering muncul di setiap pertemuan.\\\garis
Strategi soal A & Fungsi yang sudah teruji satu per satu jauh lebih mudah dirangkai daripada satu \texttt{main} panjang yang belum pernah dijalankan.\\\garis
Strategi soal B & Jawaban B tanpa langkah hanya mendapat sebagian poin, walaupun hasil akhirnya benar.\\\garis
Contoh tracing gabungan & Fungsi menghitung elemen yang lebih besar dari x secara rekursif dari belakang: 9, 9, dan 7, jadi 3.\\
\garisdasar
\end{tabular}}\par\vspace{6pt}
\Bagian{Penutup}{Cek pemahaman}
Jawab tanpa menjalankan program, lalu periksa dengan menjalankannya.
\begin{enumerate}
\item Apa yang terjadi bila elemen terakhir dari n elemen diakses dengan \texttt{a[n]}?
\item Mengapa fungsi pengurutan array of struct tidak perlu \texttt{\&} pada parameternya?
\item Bagaimana menukar dua elemen array of struct?
\item Pada rekursi, ke arah mana nilai kembali mengalir?
\item Mengapa parameter array 2D harus menyebut ukuran kolom?
\end{enumerate}
\begin{Jawaban}[]
\begin{enumerate}[leftmargin=16pt]
\item Akses di luar batas; elemen terakhir ada di \texttt{a[n - 1]}.
\item Array selalu diteruskan aslinya, bukan salinan.
\item Tukar seluruh struct, misalnya \texttt{swap(p[i], p[j])}, bukan field satu per satu.
\item Dari kasus basis naik kembali ke pemanggilan pertama.
\item Compiler perlu tahu panjang baris untuk menghitung letak \texttt{m[i][j]}.
\end{enumerate}
\end{Jawaban}
\Bagian{Penilaian}{Exit Ticket 15}
Dikerjakan di 25 menit terakhir kelas, sendiri, tanpa AI, internet, atau diskusi. Exit Ticket 15 adalah satu dari 14 Exit Ticket; rata-ratanya menjadi komponen berbobot 25\%.
\par\vspace{4pt}\noindent{\fontsize{9.5pt}{13pt}\selectfont
\begin{tabular}{@{}>{\raggedright\arraybackslash}p{0.080\textwidth}>{\raggedright\arraybackslash}p{0.600\textwidth}>{\raggedright\arraybackslash}p{0.100\textwidth}>{\raggedright\arraybackslash}p{0.200\textwidth}@{}}
\kepala{Soal} & \kepala{Isi} & \kepala{Poin} & \kepala{CPMK}\\
A1 & Tulis fungsi yang mengembalikan banyak elemen array yang lebih kecil dari rata-ratanya & 25 & CPMK0607\\\garis
A2 & Lengkapi program array of struct agar terurut berdasarkan satu field & 25 & CPMK0607\\\garis
B1 & Gambar pohon pemanggilan sebuah fungsi rekursif pada array & 20 & CPMK0608\\\garis
B2 & Jelaskan perbedaan hasil pemanggilan fungsi dengan dan tanpa referensi & 20 & CPMK0608\\\garis
B3 & Sebutkan materi pertemuan yang dipakai sebuah soal dan alasannya & 10 & CPMK0608\\
\garisdasar
\end{tabular}}\par\vspace{6pt}
\Bagian{Penutup}{Bacaan lanjut dan persiapan minggu depan}
\begin{enumerate}
\item Modul Belajar 9 sampai 14.
\item Deitel, P. dan Deitel, H. C++ How to Program, edisi ke-10, Bab 6 sampai 9.
\item Slide \texttt{01 Slide/P15 Review dan Simulasi UAS.pdf} dan hands-on di \texttt{02 Hands-on/P15 - Review UAS - Hands-on/}.
\end{enumerate}
\begin{Penting}[Minggu depan]
Pertemuan 16 adalah UAS. Kisi-kisi, tata tertib, dan contoh soal ada di slide P16 dan Modul Belajar 16.
\end{Penting}
\end{document}
